\chapter{Runtime Representation}\label{ch:runtime-representation}

Compiled code is only half of a running program.
The other half is the layout of values, objects, frames, and tables that code manipulates.
Those layouts decide the cost of every field access, call, and type test.
They also decide what the collector, the debugger, and the just-in-time compiler can observe.
This chapter fixes representations that are simple, documented, and checkable.
We start from scalar values, build up to objects and frames, and end with the maps that let tools walk the heap and the stack.

\section{Values Need Tags or Proof}\label{sec:ch32-values}

Every operation must know the kind of data it handles.
That knowledge comes from a \term{tag} stored with the value or from a proof carried elsewhere.
A tag costs bits and a test on each use.
A proof costs analysis effort and a fallback when the proof fails.
Consider adding two values in a dynamic language.
With tags, the add handler checks both operands, unboxes integers, adds, and reboxes.
With a proof, compiled code skips the checks and jumps to a slow path when a guard fails.
Figure~\ref{fig:ch32-tagged-value} shows a tagged 64-bit value with two tag bits and a 62-bit payload.
Small integers live directly in the payload.
Heap references keep a distinct tag and a pointer to an object header.
Table~\ref{tab:ch32-tagging} compares common tagging schemes by capacity, test cost, and collector support.
The table is qualitative because bit budgets depend on the target.
Pointer tagging steals low bits that alignment guarantees are zero.
Not-a-number boxing reuses quiet patterns in floating-point values.
Separate kind words avoid bit stealing at the cost of wider values.
We adopt two-bit low tags for this chapter.
Integers with tag \code{00} hold 62-bit signed payloads after an arithmetic shift.
Heap references with tag \code{01} point to eight-byte-aligned headers.
The remaining tags cover booleans with the value in the payload and a fallback box for floats and odd shapes.
Every handler states which tags it accepts before it reads any payload.

\begin{figure}[htbp]
\centering
\begin{tikzpicture}[node distance=0.6cm and 0.8cm]
\node[basic block] (tag) {tag\\ 2 bits};
\node[basic block, right=of tag, minimum width=4.5cm] (pay) {payload\\ 62 bits: integer or pointer};
\draw[cfg edge] (tag) -- (pay);
\end{tikzpicture}
\caption{A 64-bit tagged value keeps its kind in two low bits and its content in the remaining payload, so a type test is one mask and one compare.}
\label{fig:ch32-tagged-value}
\end{figure}

\begin{table}[htbp]
\centering
\caption{Tagging schemes trade value capacity against test cost and collector complexity.}
\label{tab:ch32-tagging}
\begin{tabular}{llll}
\toprule
Scheme & Capacity & Test cost & Notes \\
\midrule
Low-bit tags & 62-bit integer & Mask and compare & Needs aligned pointers \\
NaN boxing & 51-bit payload & Range compare & Shares float encoding \\
Kind word & Full width & One load & Wider values everywhere \\
Static proof & Full width & Guard on change & Needs fallback path \\
\bottomrule
\end{tabular}
\end{table}

\section{Objects, Arrays, and Headers}\label{sec:ch32-objects}

Objects add identity, layout, and method dispatch to plain values.
An \term{object header} records the kind, the size, and the information the collector needs.
A fixed-layout object stores fields at constant offsets after the header.
An array stores a length and a run of identical elements.
A hash-backed object supports changing field sets at the cost of lookup and storage overhead; whether a fixed layout is worthwhile depends on the workload.
Algorithm~\ref{alg:ch32-layout} computes field offsets for fixed-layout classes with single inheritance.
It assumes fields are listed from the root class down and each field declares a size and an alignment.
It requires that the header size and alignment are already fixed.
It returns the aligned offset of every field and the total object size.
Padding appears wherever alignment forces a gap.
Documenting each gap matters because padding is observable to native interfaces.
Method dispatch needs one more structure.
A \term{virtual table} holds one function pointer per method in a fixed order.
Each object points at the table for its class.
A call loads the table, loads the slot, and jumps indirectly.
Inline caches, covered in a later chapter, memoize the frequent target at a call site.
The layout rule stays small on purpose.
Fixed offsets make field access one load.
One header shape makes the collector walk every object with one routine.
Subclasses append fields without moving inherited offsets, so inherited code keeps working.

\begin{algorithm}[htbp]
\caption{Field offset assignment for single-inheritance classes}\label{alg:ch32-layout}
\begin{pseudocode}
\Require header size $H$ and alignment, root-to-leaf field list $F$, power-of-two alignments, checked size arithmetic
\Ensure offset $off(f)$ for each field and total size $S$
\State $cur \gets H$
\State $align_{max} \gets$ header alignment
\ForAll{$f \in F$}
\State $cur \gets$ round $cur$ up to alignment of $f$
\State $off(f) \gets cur$
\State $cur \gets cur +$ size of $f$
\State $align_{max} \gets \max(align_{max},$ alignment of $f)$
\EndFor
\State $S \gets$ round $cur$ up to $align_{max}$
\State \Return offsets and $S$
\end{pseudocode}
\end{algorithm}

Using the maximum alignment works here because power-of-two alignments divide the largest one.
For unrelated positive alignments, maximum is insufficient: the object alignment must be a common multiple, subject to the target's layout rules.
This simple runtime lays out the flattened inherited field list and permits reuse of base tail padding; a runtime that preserves the complete base size must advance past that padding before appending derived fields.

\section{Worked Example: Laying Out Two Classes}\label{sec:ch32-worked}

We lay out a root class with two fields and a subclass that adds one field.
The header occupies $16$ bytes with $8$-byte alignment.
The root declares \code{a} as a $4$-byte integer with $4$-byte alignment and \code{b} as an $8$-byte reference with $8$-byte alignment.
The subclass appends \code{c} as a $1$-byte flag with $1$-byte alignment.
We run Algorithm~\ref{alg:ch32-layout} by hand.
Initialization sets $cur = 16$ and maximum alignment to $8$.
Field \code{a} is already $4$-aligned at $16$, so $off(a) = 16$ and $cur = 20$.
Field \code{b} needs $8$-alignment, so $cur$ rounds up from $20$ to $24$, giving $off(b) = 24$ and $cur = 32$.
Field \code{c} needs only $1$-alignment, so $off(c) = 32$ and $cur = 33$.
The total rounds $33$ up to $8$-alignment, giving $S = 40$.
The four bytes from $20$ to $23$ are padding.
The seven bytes from $33$ to $39$ are trailing padding.
The listing below shows the resulting map.
\begin{irlisting}
; object layout, header 16 bytes, total 40 bytes
; off 0:  header.kind (8 bytes)
; off 8:  header.size (8 bytes)
; off 16: a (4 bytes)
; off 20: padding (4 bytes)
; off 24: b (8 bytes)
; off 32: c (1 byte)
; off 33: padding (7 bytes)
\end{irlisting}
A reader can verify each line independently.
Removing \code{c} gives total $32$ with the same offsets for \code{a} and \code{b}.
Widening \code{a} to $8$ bytes removes the interior padding and keeps total $40$.
Both variants follow from the same rounding rule, which is the property tests should pin down.

Figure~\ref{fig:visual-object-byte-map} draws the byte ranges to scale.
The one-byte field $c$ is small enough to be easily lost in a schematic box diagram, so its boundaries at $32$ and $33$ are marked separately.
Padding is storage reserved by the layout; it is not another field for the collector to scan as a reference.

\begin{figure}[htbp]
  \centering
  \begin{tikzpicture}[x=2.8mm,y=10mm]
    \filldraw[fill=black!8] (0,0) rectangle (16,1);
    \draw (16,0) rectangle (20,1);
    \filldraw[fill=black!25] (20,0) rectangle (24,1);
    \draw (24,0) rectangle (32,1);
    \filldraw[fill=black!8] (32,0) rectangle (33,1);
    \filldraw[fill=black!25] (33,0) rectangle (40,1);
    \node[font=\sffamily\small] at (8,0.5) {header};
    \node at (18,0.5) {$a$};
    \node[font=\sffamily\scriptsize] at (22,0.5) {pad};
    \node at (28,0.5) {$b$};
    \node[font=\sffamily\small] at (36.5,0.5) {padding};
    \foreach \off in {0,16,20,24,40}
      \node[font=\small,anchor=south] at (\off,1.1) {\off};
    \draw (32,1) -- (32,1.5) -- (30,1.8);
    \node[font=\small,anchor=south] at (30,1.8) {32};
    \draw (33,1) -- (33,1.5) -- (35,1.8);
    \node[font=\small,anchor=south] at (35,1.8) {33};
    \draw[cfg edge] (32.5,-0.7) -- (32.5,0);
    \node[font=\small,anchor=north] at (32.5,-0.8) {$c$: one byte};
    \node[font=\sffamily\small,anchor=north west] at (0,-0.4)
      {Byte offsets; shaded gaps contain no fields.};
  \end{tikzpicture}
  \caption{The 40-byte object contains four internal padding bytes and seven trailing padding bytes; widths are proportional to byte counts.}
  \label{fig:visual-object-byte-map}
\end{figure}



Listing~\ref{lst:visual-align-offset} supplies the checked rounding step used by the layout algorithm.
At offset $20$ with alignment $8$, the remainder is $4$, the padding is $4$, and the result is $24$.
At offset $33$, the remainder is $1$, so final rounding adds $7$ bytes and gives $40$.

\begin{listing}[htbp]
  \inputminted{c}{example-code/visual-labs/layout.c}
  \caption{Alignment rounding rejects zero alignment and checks host overflow before adding padding.}
  \label{lst:visual-align-offset}
\end{listing}

This helper accepts any positive alignment, rather than assuming a power of two.
It takes constant time under fixed-width host arithmetic and leaves its output unchanged on failure.
A complete layout implementation must also check the addition of each field's size and apply the target runtime's inheritance and alignment rules.
The test suite checks 8,192 small offset/alignment pairs against a byte-at-a-time oracle, plus zero alignment and overflow cases.

\section{Frames and Stack Maps}\label{sec:ch32-frames}

Activation records carry the same design pressure as objects.
A \term{frame} holds saved registers, locals, spill slots, and outgoing arguments in a fixed order.
The compiler emits a \term{stack map} that names every slot holding a reference at each safe point.
The collector and the deoptimizer read that map instead of guessing.
Callees address locals by frame-pointer offsets.
Callers build argument areas to a documented alignment.
A frame header records the return address, the previous frame pointer, and an index into the stack-map table.
Figure~\ref{fig:ch32-tagged-value} already showed value tags; the same discipline applies here.
Frames also contain unboxed nonreferences, return addresses, and saved machine state.
At each safepoint, the map and tag conventions must identify every live managed reference, including those held in registers, without treating arbitrary integers as pointers.
\begin{ccnote}[Maps are contracts]
A missing map entry corrupts the heap walk as surely as a wrong offset corrupts a field load.
Review map emission with the same care as code emission.
\end{ccnote}
To check maps, build a walker test with a three-deep call chain, force a collection at the deepest safepoint, and verify every live reference.
A separate deoptimization test should suspend at the same point and reconstruct locals from recovery metadata; a GC root map alone need not describe nonreference locals.
These are proposed integration tests, beyond the executable alignment helper above.

\section{Calling Conventions and Native Boundaries}\label{sec:ch32-calling}

Compiled code rarely runs alone.
It calls libraries, receives callbacks, and crosses into code the compiler never saw.
A \term{calling convention} fixes argument order, return channels, stack alignment, and caller versus callee duties.
Common conventions pass the first few arguments in registers and the rest on the stack.
The caller aligns the stack before the call and owns the outgoing area.
The callee saves the registers it disturbs and restores them before returning.
Variadic calls need extra rules because the callee cannot know the argument count from its own signature.
Native boundaries add representation duties on top of control duties.
A managed reference crossing into native code must stay visible to the collector for the whole call.
Handles or pinning scopes provide that visibility.
A native pointer entering managed code must be wrapped with a documented owner and lifetime.
Callbacks need the reverse map: a small trampoline recovers the managed closure from a native function pointer.
Test boundaries with round-trip calls that pass every supported kind in both directions.
Each test asserts the returned value, the stack alignment on entry, and that no reference was collected mid-call.
The \enquote{aligned on entry} check catches the defect where one convention assumes sixteen-byte alignment and the other guarantees only eight.

\section{Implementation Notes and Complexity}\label{sec:ch32-implementation}

Algorithm~\ref{alg:ch32-layout} runs in $O(F)$ where $F$ is the field count, using $O(F)$ space for offsets.
Field access after layout is $O(1)$.
A method call through a virtual table is $O(1)$ plus cache effects.
Stack-map size grows with the number of safe points times the live references at each point.
Sharing map encodings across nearby points keeps binaries small.
Production runtimes compress maps, split hot and cold entries, and intern common shapes.
They also align frames and headers for atomic access and vector loads.
Each refinement keeps the documented offsets stable or versions them explicitly.
A common pitfall changes a header without updating the collector, the walker, and the native interface together.
Another pitfall stores untagged raw pointers in spill slots to save one instruction.
The saved instruction costs a heap walk that cannot find the reference.
We keep every reference tagged or mapped from the start.
The layouts here feed the collector of the next chapter directly.
Headers carry the mark bits and forwarding pointers it needs.
Stack maps name the roots it must scan.
% TODO(cite): need a primary reference for stack-map encoding trade-offs.

\section{Exercises}\label{sec:ch32-exercises}

\begin{enumerate}
\item Lay out a class with fields of sizes $8$, $1$, and $4$ in that order under header $H = 16$, then reorder the fields to minimize padding and report both totals.\label{ex:ch32-layout-order}
\item Decode the bit pattern with low tag \code{01} and payload $4096$ under the two-bit scheme, state whether it is an integer or a reference, and give the denoted value or address relation.\label{ex:ch32-decode-tag}
\item Write a stack-map entry for a safe point with three live references and two integers, then describe a test that detects one wrong entry.\label{ex:ch32-map-test}
\end{enumerate}

\section{Further Reading}\label{sec:ch32-reading}

Appel covers value and frame representation as part of code generation and runtime design~\cite{appel-1998-modern}.
Morgan relates object layout and calling structure to optimization opportunities~\cite{morgan-1998-optimizing}.
Safonov examines how representation choices affect trustworthy execution~\cite{safonov-2010-trustworthy}.
The C standard draft specifies alignment, padding, and address rules for C-visible layouts~\cite{wg14-2011-c11}.
